\section{Clipping paths on arbitrary scales}\label{sec:clipping}

The quenched limit in Section~\ref{sec:gaussian} applies only along
Gaussian scales. We now obtain a weaker conclusion at every large scale:
a fixed positive amount of quenched mass has its first-hit positions in
an arbitrarily narrow tube around a suitable deterministic line. The line may depend on
the scale, but not on the environment or the starting site.

Auxiliary paths and comparison with genuine quenched mass also appear in
multiscale slowdown arguments; see Berger~\cite[Sections~5--6]{Berger2012}.
Here the retained mass is constructed from first-hit kernels using the
spatial estimates of the preceding sections.

\begin{proposition}[Clipping at arbitrary scales]\label{clip-clipping}
For any fixed \(0<T<\infty,\ \eta>0,\ p_{\rm err}>0\) there is \(d_0>0\) such that, for every sufficiently large \(r\), some deterministic slope \(s_r\) satisfies, with \(H_{\rm tot}=\lceil T n(r)\rceil\),
\[
P\left(K_{H_{\rm tot}}(x,\{\max_{j\le H_{\rm tot}}|Y_j-j s_r|\le\eta r\})\ge d_0\right)>1-p_{\rm err}.
\tag{24}\label{eq:24}
\]
The bound holds at every prescribed starting site \(x\), with the same
\(d_0\) and slope. It also gives the corresponding lower bound at every
shorter horizon by taking first-hit prefixes.
\end{proposition}

The assertion is sitewise; it does not require one high-probability event
on which all starting sites are good. This distinction permits the
averaging over initial distributions in Section~\ref{sec:kernels}.

\subsection{Transferring truncated moments}

\begin{lemma}[Truncated-moment transfer]\label{clip-moment-transfer}
Write \(\mathbb Q_h=E[K_h(0,\cdot)/a_h]\). For any \(h\to\infty,\ z\to\infty\), put \(X=W_h/z\) under \(\mathbb Q_h\), and let \(V=M_h^S/z\) have the iid sum law of \eqref{eq:8}. For fixed \(0<A<L<\infty\), with a constant independent of \(A,L\),
\[
\begin{split}
\limsup \mathbb Q_h[(X\wedge L)^2]&\le C\limsup\mathbb E[(V\wedge L)^2],\\
\limsup \mathbb Q_h[(X\wedge L)^2\mathbf1_{X>A}]&\le
 C\limsup\mathbb E[(V\wedge L)^2\mathbf1_{V>A/C}].
\end{split}
\tag{25}\label{eq:25}
\]
\end{lemma}

\begin{proof}
Take a subsequence realizing the left \(\limsup\) and extract compactified weak limits \(\mu,\nu'\) in \([0,\infty]\) for the two laws. At each \(u>0\), \(\mu([u,\infty])\) is bounded by the limit-law probability of the open event above \(u/2\), hence by the corresponding prelimit \(\liminf\). By \eqref{eq:11} and \eqref{eq:9}, this is at most \(C_1\) times a prelimit \(\limsup\) on the \(V\) side above \(u/(2C_1)\); bounding by the limiting closed tail and enlarging it once more gives
\[
\mu([u,\infty])\le C_1\nu'((u/(4C_1),\infty]).
\]
For the first inequality of \eqref{eq:25} simply integrate with weight \(2u\) up to \(L\), since capped squares are bounded continuous. For the second, first upper bound the left using \(X\ge A\) on the limiting law (upper semicontinuous test). This gives \(A^2\mu([A,\infty])\) plus the tail integral between \(A\) and \(L\). Transfer each term to bound the sum by a constant times the \(\nu'\)-expectation of the capped square on the open tail above \(A/(4C_1)\). By lower semicontinuity that last expectation is bounded by the prelimit \(\liminf\) on our subsequence. Enlarging constants gives \eqref{eq:25}.
\end{proof}

\subsection{A path policy on short blocks}

\begin{proof}[Proof of Proposition~\ref{clip-clipping}]
We will concatenate a fixed number of short blocks. In each block we
choose a probability law on a restricted set of successful prefixes,
except on a rare failure event. The two requirements are different:
the chosen endpoint displacements must have small variance, and the
restriction must retain a fixed positive amount of the original kernel
mass. The first requirement keeps the auxiliary first-hit process in a narrow tube;
the second will transfer that conclusion back to the quenched walk.

Work along an arbitrary sequence \(r\to\infty\), and pass to a
subsequence on which
\[
F(a)=\lim_{r\to\infty}\frac{m(ar)}{m(r)}
\qquad(a>0)
\]
exists. This extraction is possible because the ratios are monotone
in \(a\), and changing \(a\) by a factor changes \(m(ar)\) by at most
its square. Thus they are equicontinuous on compact logarithmic
intervals. In particular \(F(1)=1\), and the limit
\(F_0=F(0+)\) exists in \([0,1]\).
It suffices to prove the clipping assertion on every such further
subsequence with some fixed positive mass bound. Otherwise there would
be a sequence of scales at which every deterministic slope fails even
for mass thresholds tending to zero.

\paragraph{The common block rule.}
Fix a small tube width \(b>0\), and consider a block of integer height
\(h\le tn(r)\), with \(t>0\) small. Its length will be chosen
separately in the cases \(F_0=0\) and \(F_0>0\). At its current
starting site \(x\), put
\[
N_b=\{W_h\le br\},\qquad
G=K_h(x,N_b),\qquad f=G/a_h(x).
\]
Short-block survival, Equation~\eqref{eq:12}, gives a number
\(d_b\in(0,1)\), depending on the fixed \(b,t\), such that
\[
P(G<d_b)\le C_b t^2
\]
for all sufficiently large \(r\). On \(\{G<d_b\}\), mark a
failure and make an auxiliary jump to the block top with second
coordinate displacement \(b_h\) and all other horizontal
displacements zero. Otherwise sample a prefix from
\(K_h(x,\cdot\cap N_b)/G\). In the case \(F_0>0\), we will
sometimes further restrict this sampling to endpoints close to \(b_h\).

The block test and every sampled prefix use only departure rows below
the block top. Their endpoint and auxiliary choices therefore leave
fresh iid layers above the next base. Repeating any one of these
translated rules over equal blocks makes the second-coordinate endpoint
displacements \(U_h\) iid under environment averaging and auxiliary
sampling. Write \(\lambda_h=\mathbb E U_h\). Both genuine and dummy
endpoints satisfy
\[
|U_h-b_h|\le br,\qquad |\lambda_h-b_h|\le br.
\]
The dummy jumps keep the auxiliary process defined after a failure;
only trajectories without a failure will later contribute genuine
quenched path mass.

We next choose the block length and, when useful, its further endpoint
restriction. In both cases the goal is to make the variance accumulated
over \(O(n(r)/h)\) blocks small compared with \(r^2\).

\paragraph{Case \(F_0=0\): the broad restriction has small variance.}
Take \(h=\lfloor tn(r)\rfloor\) with \(t>0\) fixed and small, and
use the common block rule without further restriction. The
truncated-moment transfer and Equation~\eqref{eq:8} give
\[
\limsup_{r\to\infty}
\mathbb Q_h\!\left[\frac{W_h^2\wedge (br)^2}{r^2}\right]
\le CtF(b).
\tag{26}\label{eq:26}
\]
The normalization in the block rule costs only a fixed factor in this
moment estimate:
\[
\mathbb E[(U_h-b_h)^2]
\le4\mathbb Q_h[W_h^2\wedge(br)^2].
\]
Indeed, on \(\{f\ge1/2\}\), sampling the restriction costs at most
a factor two relative to \(K_h/a_h(x)\). On \(\{f<1/2\}\), use
\((br)^2\) and
\[
P(f<1/2)\le2\mathbb Q_h(W_h>br).
\]
Dummy jumps have zero centered displacement. Since there will be
\(O_T(1/t)\) blocks, Equation~\eqref{eq:26} makes their total
variance at most \(C_TF(b)r^2\) in the limit. This can be made as
small as needed by choosing \(b\) small.

\paragraph{Case \(F_0>0\): select a Gaussian subscale.}
Put \(z=\varepsilon r\) and \(h=\lfloor n(z)\rfloor\). Limits in
this case are taken first as \(r\to\infty\) along the chosen
subsequence, and then as \(\varepsilon\downarrow0\). We have
\[
\frac{h}{n(r)}\longrightarrow
\frac{\varepsilon^2}{F(\varepsilon)}.
\]
For every fixed \(a>0\),
\[
\limsup_{\varepsilon\downarrow0}\limsup_{r\to\infty}
 n(z)\Pr(|S|>az)=0.
\]
To see this, bound the expression by a constant depending on \(a\)
times
\((m(az)-m(az/2))/m(z)\); its iterated limit is zero because
\(F(a\varepsilon),F(a\varepsilon/2),F(\varepsilon)\to F_0>0\).

The quenched Gaussian estimate~\eqref{eq:23} consequently implies that,
for each fixed \(\rho,b>0\), there is \(\alpha_\rho>0\) such that
\[
\limsup_{\varepsilon\downarrow0}\limsup_{r\to\infty}
P\left((K_h/a_h)\{W_h\le br,\ |Y_h-b_h|\le\rho z\}
       <\alpha_\rho\right)=0.
\tag{27}\label{eq:27}
\]
Here and below a block kernel without a stated start is at the origin.
For completeness, a failure of this implication would give
\(\varepsilon_j\downarrow0\) and sufficiently large \(r_j\) for
which the first \(j\) large-jump tests at thresholds
\(a=1,1/2,\ldots,1/j\) are less than \(1/j\), whereas the
probability in \eqref{eq:27} stays bounded away from zero. Taking
\(z_j=\varepsilon_jr_j\to\infty\) gives a Gaussian sequence, and
\(b/\varepsilon_j\to\infty\) is an allowed tube multiple in
\eqref{eq:23}. This is a contradiction.

We also need a moment estimate for the fallback to the broad
restriction:
\[
\limsup_{\varepsilon\downarrow0}\limsup_{r\to\infty}
\frac{\mathbb Q_h[(W_h^2\wedge(br)^2)
                         \mathbf1_{\{W_h>Az\}}]}{z^2}
\le \frac{C}{A^2}+C\frac{F(b)-F_0}{F_0}.
\tag{28}\label{eq:28}
\]
Fix \(A,b,\varepsilon\) with \(b>A\varepsilon\) and
\(b>\varepsilon\), and apply Equation~\eqref{eq:25} with
\(L=b/\varepsilon\). For the resulting iid sums, split each
symmetric increment into its parts on
\[
\{|S|\le z\},\qquad
\{z<|S|\le br\},\qquad
\{|S|>br\}.
\]
All three parts are centered. The maximal sum \(U\) of the first
parts has fourth moment at most \(Cz^4\), since \(hm(z)\le z^2\).
The second maximal sum \(U'\) has second moment bounded by
\(Ch\mathbb E[S^2;z<|S|\le br]\). The probability that any third
part occurs, multiplied by the cap \((br)^2\), is bounded by
\(h(br)^2\Pr(|S|>br)\). These last two bounds add to at most
\[
Ch\bigl[m(br)-m(z)+z^2\Pr(|S|>z)\bigr].
\]
Without a third jump, the original maximal sum is at most \(U+U'\).
On the event that it exceeds \(Az\), its capped square is bounded
by a fixed constant times
\(U^2\mathbf1_{\{U>Az/C'\}}+(U')^2\).
Divide by \(z^2\), use the fourth-moment bound for the first term,
and take the indicated iterated limits. The last large-jump term
vanishes; this proves \eqref{eq:28}.

We now refine the common block rule. Let
\[
g=(K_h/a_h(x))\{W_h\le br,\ |Y_h-b_h|\le\rho z\}.
\]
On a nonfailure, when \(g\ge\alpha_\rho\), sample this stricter
restriction, normalized. When \(g<\alpha_\rho\), use the broad
restriction as before. Take \(\alpha_\rho\le1\). The raw mass of
whichever restriction is selected is then at least
\(\alpha_\rho d_b\): in the stricter case this follows from
\(a_h(x)\ge G\ge d_b\).

For \(Az<br\), the endpoint displacement of this refined rule
satisfies
\[
\frac{\mathbb E[(U_h-b_h)^2]}{z^2}
\le \rho^2+A^2P(g<\alpha_\rho)
 +4\frac{\mathbb Q_h[(W_h^2\wedge(br)^2)
                     \mathbf1_{\{W_h>Az\}}]}{z^2}.
\tag{29}\label{eq:29}
\]
The stricter restriction contributes at most \(\rho^2\).
In the fallback case with \(f\ge1/2\), split at \(W_h=Az\) and
use the factor-two normalization bound. When \(f<1/2\), use the
same bound on its probability as in the first case, noting that
\(\{W_h>br\}\subset\{W_h>Az\}\). Again the dummy displacement
costs zero. Equations~\eqref{eq:27}--\eqref{eq:28} show that the
iterated limit superior in \eqref{eq:29} can be made arbitrarily small:
first choose \(\rho,b\) small and \(A\) large, then choose
\(\varepsilon\) small.

\paragraph{Concatenation and transfer to genuine path mass.}
In either case, run \(J=\lceil H_{\rm tot}/h\rceil\) blocks with
the selected rule. In the second case set
\(t=2\varepsilon^2/F_0\), which bounds \(h/n(r)\) for large \(r\).
For small fixed \(\varepsilon\), also
\(F(\varepsilon)\le2F_0\), so in both cases
\(J=O_T(1/t)\). The probability of any marked failure is at most
\[
J C_bt^2=O_{b,T}(t).
\]
For the centered sums of the iid endpoint displacements, the maximal inequality for independent centered sums gives
\[
\Pr\left(\max_{j\le J}
 \left|\sum_{u=1}^j(U_h^{(u)}-\lambda_h)\right|
                       \ge\eta r/2\right)
\le \frac{4J\mathbb E[(U_h-b_h)^2]}{\eta^2r^2}.
\]
In the first case the limit superior of the numerator divided by
\(r^2\) is at most \(C_TF(b)\). In the second case
\(Jz^2/r^2\le C_T\), and Equation~\eqref{eq:29} gives the same
required smallness. Choose the parameters in the orders stated above. In the first case
choose \(t\) sufficiently small after \(b\); in the second choose
\(\varepsilon\) sufficiently small after \(\rho,b,A\), with
\(t=2\varepsilon^2/F_0\). In both cases require the failure probability
to be small and
\[
2b+C\sqrt t<\eta/2,
\]
where \(C\) is the constant in Equation~\eqref{eq:10}.

Within a genuine block, its median-tube restriction and
Equation~\eqref{eq:10} imply
\[
\max_{j\le h}|Y_j-j\lambda_h/h|
 \le (2b+C\sqrt t)r<\eta r/2.
\]
Thus, outside the failure and maximal-deviation events, the concatenated
first-hit positions through \(Jh\) stay within \(\eta r\) of the
deterministic line of slope \(s_r=\lambda_h/h\).

Choose the two annealed error bounds small enough that, with environment
probability greater than \(1-p_{\rm err}\), the conditional auxiliary
law at the specified start assigns at least one half to these genuine
tube paths. On a nonfailure, its density with respect to the local
kernel is at most \(d_*^{-1}\), where
\(d_*=d_b\) in the first case and
\(d_*=\alpha_\rho d_b\) in the second. The density of a genuine
concatenation is therefore at most \(d_*^{-J}\).
The concatenated local kernels are submeasures of the original-floor
kernel, by the quenched Markov property at first hits: each local
no-drop condition merely adds a restriction. All parameters are now
fixed, so \(J\le J_*<\infty\) for large \(r\). The original kernel
through \(Jh\) consequently gives the tube mass at least
\(d_*^{J_*}/2\). Taking its first-hit prefix at \(H_{\rm tot}\)
proves Equation~\eqref{eq:24}.
\end{proof}
